\documentclass[11pt,letterpaper]{article}
\usepackage[utf8]{inputenc}
\usepackage[T1]{fontenc}
\usepackage[spanish,es-nodecimaldot,es-tabla]{babel}
\usepackage{mathptmx}
\usepackage{amsmath,amssymb}
\usepackage{graphicx}
\usepackage{booktabs,array}
\usepackage{tabularx}
\usepackage{float}
\usepackage{enumitem}
\usepackage[hyphens]{url}
\usepackage[hidelinks]{hyperref}
\usepackage{geometry}
\usepackage{xcolor}
\usepackage[most]{tcolorbox}
\usepackage{pdflscape}
\geometry{letterpaper,left=2.4cm,right=2.4cm,top=2.3cm,bottom=2.3cm}
\linespread{1.1}
\setlength{\parindent}{0pt}
\setlength{\parskip}{0.5em}
\setlist[itemize]{leftmargin=1.4em,topsep=1pt,itemsep=1pt}
\setlist[enumerate]{leftmargin=1.7em,topsep=1pt,itemsep=1pt}
\definecolor{iiBlue}{HTML}{1F4E79}
\definecolor{iiOrange}{HTML}{C55A11}
\newtcolorbox{propbox}[1]{colback=iiBlue!6,colframe=iiBlue,coltitle=white,fonttitle=\bfseries,title=#1,arc=2pt,boxrule=0.9pt,left=6pt,right=6pt,top=3pt,bottom=3pt}
\newtcolorbox{alertbox}[1]{colback=iiOrange!7,colframe=iiOrange,coltitle=white,fonttitle=\bfseries,title=#1,arc=2pt,boxrule=0.9pt,left=6pt,right=6pt,top=3pt,bottom=3pt}
\newcommand{\hrulethin}{\noindent\rule{\textwidth}{0.6pt}}
\begin{document}\pagestyle{plain}
\begin{center}
    \textbf{\large Universidad de Santiago de Chile}\\[0.1cm]
    Departamento de Ingeniería Industrial
\end{center}
\vspace{0.15cm}\hrulethin\vspace{0.35cm}
\begin{center}
    {\scshape Estadística Aplicada \quad$\cdot$\quad Caso 1}\\[0.35cm]
    {\LARGE \textbf{Enunciado}}\\[0.15cm]
    {\large Proyecto Integrador: Regresión}\\[0.2cm]
    Predicción de una respuesta continua
\end{center}
\vspace{0.2cm}\hrulethin\vspace{0.4cm}
\begin{propbox}{Propósito}
Definir desde cero un problema de \textbf{predicción de una respuesta continua} en Ingeniería Industrial y
construir un sistema reproducible que integre las tres unidades del curso: los \textbf{fundamentos} (decisión,
población, incertidumbre e inferencia), la \textbf{regularización} (lineal, ridge, lasso, elastic-net,
splines/GAM y validación cruzada) y los \textbf{ensembles} (árbol, Random Forest, boosting y combinación).
\end{propbox}

\textbf{Dinámica del curso.} El curso contempla dos casos de estudio. El \textbf{Caso 1 (Regresión)} es común a todos los equipos. En el \textbf{Caso 2}, cada equipo recibe aleatoriamente una variante: \textbf{clasificación supervisada} o \textbf{clasificación no supervisada}. Este documento corresponde a el \textbf{Caso 1: Regresión}.

\section*{1. Problema y dataset}
Cada equipo selecciona un problema realista de Ingeniería Industrial (demanda, consumo energético, tiempos de proceso, costos, calidad dimensional, productividad, degradación o mermas) y busca un dataset externo. Antes de modelar, entrega una propuesta de una página con: contexto, usuario y decisión apoyada; unidad de análisis (entidad y momento), respuesta continua, instante de predicción, horizonte y población objetivo; fuente, enlace, licencia y responsable; tamaño, variables, período y criterios de inclusión; y riesgos iniciales (leakage, sesgo, dependencia, calidad). Mínimo sugerido: 1.000 observaciones utilizables y al menos ocho variables informativas. No se permiten datasets usados en clases.

\section*{2. Métodos mínimos}
\begin{enumerate}
\item baseline (media, valor previo o promedio por grupo)
\item regresión lineal (OLS) con diagnóstico de residuos y multicolinealidad
\item ridge, lasso y elastic-net
\item modelo flexible: splines o GAM
\item árbol de regresión y Random Forest
\item Gradient Boosting o XGBoost
\item combinación de modelos (averaging o stacking)
\end{enumerate}
Cada exclusión debe justificarse con evidencia (error fuera de muestra, estabilidad o interpretabilidad).

\section*{3. Pregunta de decisión}
¿Qué modelo predice mejor la respuesta continua fuera de muestra, con qué error e incertidumbre, y qué
complejidad se justifica frente a un baseline?

\section*{4. Protocolo}
\begin{itemize}
\item Diseñar una partición coherente con el uso futuro (temporal, por grupos o aleatoria).
\item Estimar al menos una relación con intervalo de confianza y distinguir significancia de relevancia.
\item Seleccionar hiperparámetros ($\lambda$, $\alpha$, nudos, \texttt{mtry}, profundidad, \texttt{shrinkage}) por validación cruzada.
\item Ajustar imputación, codificación, escalamiento y selección dentro de cada fold; fijar semillas.
\item Comparar siempre frente a un baseline o solución de referencia, con incertidumbre.
\item Congelar el sistema antes de abrir el test y evaluar al menos un subgrupo relevante.
\end{itemize}
\begin{alertbox}{Integridad}
El uso anticipado del test, el preprocesamiento global o la selección con test invalidan la evaluación de generalización. Una buena métrica no compensa fuga de información ni la ausencia de baseline.
\end{alertbox}

\section*{5. Equipo}
Trabajo en equipo. No se asignan roles fijos: todos los integrantes deben participar en la presentación y poder defender cualquier parte (datos, código, método y decisiones).

\section*{6. Entregables y formato}
\begin{verbatim}
equipo_regresion/
|-- presentacion.pdf   (o .pptx)   <- producto principal
|-- codigo/            (script reproducible en R)   <- producto principal
|-- informe.pdf        (respaldo, maximo 5 paginas)
`-- datos/ o enlace_descarga.txt
\end{verbatim}
El \textbf{producto evaluado es la presentación y el script}; el informe es un respaldo (máximo 5 páginas, sin
contar portada ni referencias). La presentación tiene \textbf{10 diapositivas, 10 minutos de exposición y 10
minutos de preguntas}. El script debe ejecutarse desde una sesión limpia y reproducir los resultados. Un
diccionario de variables se incluye como anexo del informe.

\section*{7. Evaluación}
La evaluación se concentra en el \textbf{producto}: la calidad de la presentación (diseño y oratoria) y del
script. El informe respalda las afirmaciones, pero no es el foco. La mayor ponderación recae en la
\textbf{oratoria y dominio conceptual} y en la \textbf{calidad de la presentación}. Se califica con \textbf{dos notas independientes}: \textbf{cátedra} (presentación, oratoria, rigor e informe de respaldo) y \textbf{laboratorio} (el script). La ponderación detallada está en la rúbrica. Existe además una \textbf{co-evaluación} entre pares que puede reducir las notas individuales (ver rúbrica).

\section*{8. Tabla obligatoria (validación)}
\begin{center}\small
\begin{tabular}{lllllll}
\toprule
Modelo & Familia & RMSE & MAE & $R^2$ & RMSE subgrupo & vs baseline \\
\midrule
Baseline & --- & & & & & \\
OLS & Lineal & & & & & \\
Ridge / Lasso / EN & Regularizado & & & & & \\
Splines / GAM & Flexible & & & & & \\
Random Forest & Ensemble & & & & & \\
Boosting & Ensemble & & & & & \\
Stacking / averaging & Ensemble & & & & & \\
\bottomrule
\end{tabular}\end{center}

El resultado final de test se presenta por separado, sólo para el sistema seleccionado, frente al baseline. Indique ``n/a'' donde una métrica no aplique.

\section*{9. Hitos}
\begin{itemize}
\item \textbf{Semana 1:} propuesta aprobada; datos, auditoría, partición y baseline.
\item \textbf{Semana 2:} modelos, validación y comparación preliminar.
\item \textbf{Semana 3:} test, presentación, script e informe de respaldo.
\end{itemize}\end{document}
